\section{General Relativistic Stellar Structure and the TOV Equation}
\label{sec:tov}

The Newtonian analysis of Section~\ref{sec:newtonian} identified the
correct physical mechanism underlying the neutron star maximum mass but
overestimated it by nearly an order of magnitude. Two approximations must
be relaxed to obtain a quantitatively reliable result: the equation of
state must be taken in its full relativistic form, and the gravitational
field must be treated within general relativity rather than Newtonian
theory. The necessity of the latter is quantified by the compactness
parameter,
%
\begin{equation}
    \mathcal{C} \equiv \frac{GM}{Rc^2}
                \sim \frac{G \cdot 1.4\,M_\odot}{(10~\text{km})\,c^2}
                \sim 0.2,
    \label{eq:compactness}
\end{equation}
%
which measures the ratio of the Schwarzschild radius $r_s = 2GM/c^2$ to
the stellar diameter. For the Sun $\mathcal{C} \sim 10^{-6}$, and
Newtonian gravity is an excellent approximation. For a neutron star
$\mathcal{C} \sim 0.2$, placing it firmly in the regime where
general relativistic corrections are of order 20\% and cannot be
neglected. This section derives the Tolman--Oppenheimer--Volkoff (TOV)
equation, the general relativistic generalization of hydrostatic
equilibrium, and uses it together with the equation of state of
Section~\ref{sec:fermi} to determine the maximum stable neutron star mass.

% ------------------------------------------------------------
\subsection{The Schwarzschild Metric and Stellar Structure}
% ------------------------------------------------------------

For a static, spherically symmetric mass distribution the spacetime
geometry is described by the Schwarzschild-like interior metric,
%
\begin{equation}
    ds^2 = -e^{2\Phi(r)}c^2\,dt^2
           + \left(1 - \frac{2Gm(r)}{rc^2}\right)^{-1}dr^2
           + r^2\,d\Omega^2,
    \label{eq:interior-metric}
\end{equation}
%
where $\Phi(r)$ is the gravitational potential and $m(r)$ is the
gravitational mass enclosed within radius $r$ \cite{Misner1973}. The
mass function $m(r)$ satisfies the same continuity equation as in the
Newtonian case,
%
\begin{equation}
    \frac{dm}{dr} = 4\pi r^2 \varepsilon(r)/c^2,
    \label{eq:mass-continuity-gr}
\end{equation}
%
except that the source term is now the total energy density $\varepsilon$
(including rest mass) rather than the rest mass density $\rho$ alone,
reflecting the fact that in general relativity all forms of energy
gravitate.

\subsection{Derivation of the TOV Equation}
 
The equation of hydrostatic equilibrium in general relativity is obtained
from the covariant conservation of the stress-energy tensor,
$\nabla_\mu T^{\mu\nu} = 0$, applied to a perfect fluid with energy
density $\varepsilon$ and isotropic pressure $P$ \cite{Tolman1939,
Oppy}. The $r$-component of this conservation law, evaluated
using the metric \eqref{eq:interior-metric}, yields the
Tolman--Oppenheimer--Volkoff equation,
%
\begin{equation}
    \frac{dP}{dr} = -\frac{G}{r^2}
    \left(\varepsilon + \frac{P}{c^2}\right)
    \left(m(r) + \frac{4\pi r^3 P}{c^2}\right)
    \left(1 - \frac{2Gm(r)}{rc^2}\right)^{-1},
    \label{eq:tov}
\end{equation}
%
which recovers the Newtonian equation~\eqref{eq:hydrostatic} in the
limit $c \to \infty$. The TOV equation, the mass continuity
equation~\eqref{eq:mass-continuity-gr}, and the equation of state
$P = P(\varepsilon)$ from Section~\ref{sec:fermi} form a closed system
of ordinary differential equations in the fully general relativistic
setting.

\subsection{Numerical Integration and Boundary Conditions}
 
The TOV system is integrated numerically outward from $r = 0$ with
boundary conditions $m(0) = 0$ and $P(0) = P(\varepsilon_c)$, where
$\varepsilon_c$ is the chosen central energy density. Integration
proceeds until $P(R) = 0$, which defines the stellar surface at radius
$R$, and the total gravitational mass is $M = m(R)$. The equation of
state is evaluated numerically from
equations~\eqref{eq:energy-density-rel} and~\eqref{eq:pressure-rel},
with the Fermi momentum $p_F$ as the parameter relating $\varepsilon$
and $P$ implicitly through equation~\eqref{eq:number-density}.
Repeating the integration for a sequence of central densities
$\varepsilon_c$ traces out the full mass-radius curve.


% ------------------------------------------------------------
\subsection{The Oppenheimer--Volkoff Maximum Mass}
% ------------------------------------------------------------

The central result of the TOV integration with the free neutron Fermi
gas equation of state is the existence of a maximum stable mass. As the
central density increases, the mass $M(\varepsilon_c)$ rises to a
maximum and then decreases; configurations on the decreasing branch are
unstable to radial perturbations \cite{shapiro_teukolsky_1983}. The maximum mass
obtained from this calculation is the Oppenheimer-Volkoff limit,
%
\begin{equation}
    M_{\mathrm{OV}} \approx 0.7\, M_\odot,
    \label{eq:ov-limit}
\end{equation}
%
occurring at a central density $\varepsilon_c \approx 5 \times 10^{15}$~g~cm$^{-3}$
and a corresponding radius $R \approx 9.6$~km \cite{Oppy}.
This is the neutron star analog of the Chandrasekhar limit for white
dwarfs, and it marks the boundary between stable neutron stars and
objects that must collapse further --- in the absence of additional
pressure support --- to black holes.

Comparing equations~\eqref{eq:chandrasekhar-neutron} and
\eqref{eq:ov-limit}, the inclusion of general relativity reduces the
maximum mass from the Newtonian estimate of $\sim 5.4\,M_\odot$ to
$0.7\,M_\odot$, a reduction by nearly an order of magnitude. This
dramatic difference reflects the cumulative effect of the three GR
correction terms identified in Section~5.3: each individually increases
the effective gravitational pull, and together they render the star
unstable at a far lower mass than Newtonian gravity would predict. The
physical content of the OV limit is therefore not just a quantitative
correction to the Chandrasekhar result but a qualitatively different
regime in which spacetime curvature plays a central role in determining
stellar stability.